\subsection{Einstein-chain theorem and falsification boundary}

The finite support and algebraic packets establish repair-invariant incidence,
the spherical branch receipts, cap-normal identities, the Lorentz action on
the celestial screen, exact finite entropy splitting on the central-interface
branch, and the finite tomography identities used to reconstruct a trace-free
symmetric tensor from null projections.

Authenticated semantic read-from provenance supplies the finite event carrier
and its generated precedence order. Its canonical source height is the
attained longest authenticated-parent-chain length and strictly increases on
generated ancestry.
Separately, every finite decidable strict partial order compiles into an
abstract authenticated log whose direct parent relation and generated order
are exactly the supplied relation. Thus the generic parent relation contains
every supplied predecessor rather than Hasse covers alone; this proves grammar expressivity, not
dynamical or physical selection.
Independently, the exact rank-three positive source Gram quotient and a real
axis define the four-dimensional ambient target carrier
\[
 \mathbb R\oplus V_{\rm src},\qquad
 Q(t,x)=t^2-g_{\rm src}(x,x),
\]
with Lorentz inertia \((1,3)\); source-unit directions are exactly its
future-null rays. Canonical source height enters only as the temporal
coordinate of a positively scaled event placement. Every finite log has an
explicit auxiliary injective forward-causal placement obtained from a finite
enumeration along one source axis and a scale larger than its diameter. It is
not order-reflecting or source-selected. A supplied physical-candidate spatial
readback and edge speed bound give
one-way cone compatibility. Equal-height spatial separation and strict
spacelikeness of unsupported increasing-height pairs imply converse support
and exact two-way order--cone equivalence. Antisymmetry derives injectivity;
the two-way equivalence gives exact interval preservation. Physical event identity,
the spatial readback, the order-faithful-placement conditions, and operational
clock calibration remain separately testable.
A controlled effective route uses vanishing causal discrepancy and
compatible count-volume convergence. Conservative Fibonacci-record
populations with a specified complete-neighbour read law have such a flat
$1+3$ limit when their fill error is small relative to the shrinking read
radius. Equal-mass assignments with vanishing displacement error give
normalized event counts converging to volume. On fixed interior timelike
diamonds, fourth roots of interval-count ratios recover proper-time ratios,
with a reference interval fixing the
unit and complete ancestry access supplying the readout. Within the class of
nonzero closed convex pointed displacement cones, operational invariance
under source rotations and every boost along one axis forces the Lorentz
cone up to time orientation. The covariance is an assumption about actual
influences and readouts. Physical selection of the population, read law and
common matter realization is separate from these conditional results.

The source-direction Einstein theorem works in a separately supplied \(3+1\)
tensor interface indexed by the same finite event type and takes balance on nine
fixed algebraic inverse images of the coordinate null-tomography directions,
symmetric geometry/stress fields in that tensor interface, discrete Ward and
Bianchi identities, an algebraic coupling, four step maps, a base event, and
connectedness. It
derives all-null balance, the pointwise and constant metric ambiguities, and
the Einstein-form tensor identity in that supplied tensor interface. The nine directions are
not observed provenance links, physical rays, or sky samples, and the supplied
fields are not thereby identified with curvature or physical stress-energy.
The informational order shares the event type but is not used in the tensor
calculation: no provenance link selects a direction, field, step, or balance.
The distinguished nine-set is not selected by the source dynamics and is not
invariant as a set under Lorentz or \(\mathrm{SO}(3)\) transformations. The
matrix fields and discrete step maps remain supplied tensor data, and their
coordinate differences are not identified with those of the constructed
carrier. Ambient target-carrier dimension and signature are constructed
independently by the source-causal carrier theorem; intrinsic poset dimension
remains a continuum diagnostic. Neither those fields nor the physical premise
bundle is constructed.

On the exact-embedding route, an effective Lorentzian manifold requires the source-causal
continuum certificate of Theorem 4.3e: one physical refinement family with
exact two-way order/cone agreement, calibrated count--volume, independent
dimension and manifoldlikeness tests, stable topology, and a unique
distinguishing Lorentzian limit. The smooth Einstein reading additionally
requires a same-family tensor-curvature reconstruction converging to the
smooth Einstein tensor, or the separately stated continuum
small-ball/null-balance identification, plus modular and scaling control,
physical stress, Ward/Bianchi, universal coupling, a vacuum reference, and
physical \(8\pi G\) scale identification on that same family. A convergent
causal-set scalar-curvature estimator is an independent diagnostic and is
insufficient by itself. Finite certificate programs verify
their stated algebra and controls; no source-derived family carrying the full
continuum antecedent is constructed here.

The implication requires its receipts on one cofinal branch, a rank-three
source Gram quotient, tensor-linearity of the directional charges, the
stated curvature or small-diamond remainder rate, and scale readouts with
trivial common rescaling. The chosen causal route must meet its exact
order--cone criterion or its controlled causal and volume error bounds.

\begin{center}\rule{0.5\linewidth}{0.5pt}\end{center}
